\section{الگوریتم‌های بالاآوردن}

پاندول معکوس با چرخ واکنشی یک سیستم غیرخطی و جفت‌شده است که در آن، چرخ واکنشی از طریق موتور جریان مستقیم (اتصال مستقیم، $N=1$) به پاندول کوپل شده است. هدف کنترلرهای بالاآوردن (Swing-Up) رساندن پاندول از وضعیت آویزان به وضعیت عمودی است. در این بخش، سه روش بالاآوردن پیاده‌سازی‌شده را بررسی می‌کنیم و سپس نحوه انتقال کنترل به کنترلرهای تعادلی را توضیح می‌دهیم.

برای مدل کامل دینامیک سیستم، شامل ماتریس اینرسی کوپل‌شده و معادلات مدار آرمیچر، به بخش~\ref{sec:system-math} مراجعه شود.

\subsection{مدل سیستم و مرجع انرژی}

پاندول و چرخ واکنشی یک سیستم دورانی دو‌جسمی تشکیل می‌دهند. اینرسی مؤثر چرخ، شامل اینرسی روتور موتور (در اتصال مستقیم $N=1$)، به صورت زیر تعریف می‌شود:

\begin{equation}
    I_w^{\text{eff}} = I_w + I_{\text{rotor}} \, N^2
\end{equation}

معادلات حرکت کوپل‌شده به شکل ماتریسی عبارتند از:

\begin{equation}
    \begin{bmatrix} M_{11} & M_{12} \\ M_{12} & M_{22} \end{bmatrix}
    \begin{bmatrix} \ddot{\theta} \\ \ddot{\varphi} \end{bmatrix}
    =
    \begin{bmatrix} G \sin\theta - N\,K_t\, i_a - b\,\dot\theta \\ N\,K_t\, i_a - b_w^{\text{eff}}\,\dot\varphi \end{bmatrix}
\end{equation}

که در آن جملات اینرسی به صورت زیر تعریف می‌شوند:

\begin{align}
    M_{11} &= I_p + m_w L^2 + I_w^{\text{eff}} \\
    M_{12} &= M_{22} = I_w^{\text{eff}}
\end{align}

و ضریب گرانشی، مجموع اثر وزن پاندول و چرخ را در بازوهای گشتاور مربوطه تجمع می‌دهد:

\begin{equation}
    G = \left(m_p \, \ell_{\text{com}} + m_w \, L\right) g
\end{equation}

در روابط بالا، $I_p$ اینرسی پاندول حول محور چرخش، $m_p$ و $m_w$ جرم پاندول و چرخ، $\ell_{\text{com}}$ فاصله مرکز جرم پاندول از محور، $L$ طول پاندول و $g$ شتاب گرانش است.

انرژی مکانیکی پاندول نسبت به پیکربندی سکون در وضعیت عمودی ($\theta = 0$، $\dot\theta = 0$) مرجع‌دهی می‌شود:

\begin{equation}
    E_p(\theta, \dot\theta) = \tfrac{1}{2}\, I_p\, \dot\theta^2 + G\,(\cos\theta - 1)
\end{equation}

در وضعیت عمودی، $E_p = 0$ و در پایین‌تر از عمودی، $E_p < 0$ است. هدف بالاگردان، رساندن هم‌زمان $E_p \to 0$ و $\dot\theta \to 0$ است.

\subsection{بالاآوردن مبتنی بر انرژی (روش آستروم–فوروتو)}

\subsubsection{شهود فیزیکی}

پاندول یک نوسانگر غیرخطی است. با اعمال گشتاور چرخ در فاز صحیح با نوسان پاندول، در هر نیم‌سیکل انرژی به سیستم تزریق می‌شود؛ دقیقاً مشابه عمل کودکی که با جابه‌جایی مرکز جرم خود در زمان مناسب، تاب را به ارتفاع بالاتر می‌رساند. نکته کلیدی روش آستروم–فوروتو (Åström–Furuta) این است که فقط انرژی پاندول ردیابی می‌شود، نه انرژی کل سیستم. به این ترتیب، انرژی به‌صورت غیرمفید در چرخ در حال چرخش انباشته نمی‌شود.

\subsubsection{قانون کنترل}

فرمان ولتاژ متناسب با حاصل‌ضرب کسری انرژی و سرعت چرخ است:

\begin{equation}
    V = -k_e \;\bigl(E_{\text{target}} - E_p\bigr)\; \dot\varphi
\end{equation}

که در آن $k_e > 0$ بهره انرژی (پارامتر \texttt{energy\_swing\_up\_gain}، مقدار پیش‌فرض $1.0$)، $E_{\text{target}} = 0$ ژول (انرژی سکون عمودی) و $\dot\varphi$ سرعت زاویه‌ای نسبی چرخ است. علامت منفی تضمین می‌کند که وقتی انرژی کمتر از هدف است، ولتاژ چرخ را در جهتی می‌راند که انرژی پاندول در نیم‌نوسان فعلی افزایش یابد.

\subsubsection{میرایی انرژی اضافی}

اگر پاندول از هدف انرژی فراتر رود ($E_p > 0$)، ادامه تزریق انرژی باعث چرخش پیوسته پاندول می‌شود. در این حالت، کنترلر به میرایی چرخ تغییر وضعیت می‌دهد:

\begin{equation}
    V = -\tfrac{1}{2}\, k_e\, \dot\varphi
\end{equation}

این فرمان، انرژی جنبشی چرخ را از طریق تلفات مقاومتی موتور تخلیه می‌کند.

\subsubsection{تحریک فازى}

در نزدیکی نقاط بازگشت نوسان، سرعت چرخ $\dot\varphi$ به صفر نزدیک می‌شود و سیگنال تزریق انرژی $-k_e (E_{\text{target}} - E_p) \dot\varphi$ از بین می‌رود. در نتیجه، پاندول متوقف می‌شود. برای فرار از این بن‌بست، هرگاه $|\dot\varphi| < \varphi_{\text{exc}}$ باشد (که $\varphi_{\text{exc}} = 0.5$ رادیان بر ثانیه است)، یک ولتاژ تحریک تزریق می‌شود:

\begin{equation}
    V_{\text{exc}} = \alpha\, V_{\max}\, \operatorname{sign}(\sin\theta)
\end{equation}

که در آن $\alpha = 0.3$ کسر تحریک از ولتاژ بیشینه است. جهت $\operatorname{sign}(\sin\theta)$ تضمین می‌کند که ضربه چرخ در جهتی باشد که پاندول را در نوسان بعدی به سمت بالا ببرد.

\subsubsection{کاهش تدریجی سرعت چرخ}

برای جلوگیری از شتاب‌گیری بیش از حد چرخ، ولتاژ به‌صورت خطی کاهش می‌یابد وقتی $|\dot\varphi|$ از $60\%$ سرعت بیشینه مجاز چرخ $\omega_{\max}$ فراتر رود:

\begin{equation}
    \text{scale} = \max\!\left(0,\; \frac{\omega_{\max} - |\dot\varphi|}{\omega_{\max} - 0.6\,\omega_{\max}}\right), \qquad V \leftarrow V \cdot \text{scale}
\end{equation}

این رابطه یک شیب خطی از ولتاژ کامل در $|\dot\varphi| = 0.6\,\omega_{\max}$ تا صفر در $|\dot\varphi| = \omega_{\max}$ ایجاد می‌کند.

\subsection{بالاآوردن با خطی‌سازی بازخورد جزئی}

\subsubsection{شهود فیزیکی}

روش خطی‌سازی جزئی بازخورد (Partial Feedback Linearization) رویکردی متفاوت از تزریق انرژی دارد. در این روش، به‌جای اینکه به‌صورت غیرمستقیم انرژی پاندول را افزایش دهیم، شتاب زاویه‌ای پاندول را مستقیماً فرمان می‌دهیم. ایده اصلی این است: اگر بتوانیم مشخص کنیم پاندول در هر لحظه چه شتابی باید داشته باشد، می‌توانیم از طریق معکوس‌سازی دینامیک کوپل‌شده، شتاب چرخ مورد نیاز برای تولید آن شتاب پاندول را محاسبه کنیم و سپس این شتاب چرخ را به ولتاژ موتور تبدیل نماییم.

این روش مبتنی بر مدل است؛ یعنی به دانش دقیق ماتریس اینرسی $\mathbf{M}$، ضریب گرانشی $G$ و ضرایب میرایی نیاز دارد. در مقابل روش انرژی که فقط $E_p$ را ردیابی می‌کند، خطی‌سازی بازخورد جزئی از ساختار کامل معادلات حرکت بهره می‌گیرد.

\subsubsection{شتاب مطلوب پاندول}

نخستین گام، تعریف یک شتاب مطلوب برای پاندول است. این شتاب، رفتار موردنظر پاندول را بدون توجه به چگونگی تولید آن مشخص می‌کند. برای این منظور، از یک قانون مجازی \LR{PD} (تناسبی--مشتقی) استفاده می‌کنیم. ایده این قانون ساده است: بخش تناسبی، شتابی متناسب با انحراف زاویه‌ای پاندول از وضعیت عمودی تولید می‌کند (مانند یک فنر مجازی) و بخش مشتقی، شتابی متناسب با سرعت زاویه‌ای پاندول در جهت مخالف حرکت اعمال می‌نماید (مانند یک میراگر مجازی). بدون بخش مشتقی، پاندول حول نقطه هدف نوسان بی‌پایان انجام می‌داد؛ این بخش، انرژی جنبشی مجازی را در هر نیم‌سیکل تخلیه می‌کند و همگرایی مجانبی را تضمین می‌نماید.

قانون مجازی \LR{PD} به صورت زیر تعریف می‌شود:

\begin{equation}
    \ddot\theta_{\text{des}} = -k_p \sin\theta - k_d\, \dot\theta
\end{equation}

در این رابطه، $k_p^{\text{PFL}}$ بهره تناسبی و $k_d^{\text{PFL}}$ بهره مشتقی قانون مجازی \LR{PD} هستند.

\paragraph{نقش فیزیکی $k_p^{\text{PFL}}$ (بهره تناسبی).}
جمله $-k_p \sin\theta$ یک شتاب بازگرداننده وابسته به موقعیت ایجاد می‌کند؛ معادل یک فنر مجازی غیرخطی که پاندول را به سمت $\theta = 0$ می‌کشاند. استفاده از $\sin\theta$ به‌جای $\theta$ ضروری است، زیرا این کنترلر باید در تمام زوایا (نه فقط نزدیکی عمودی) کار کند. در نزدیکی افقی، $|\sin\theta| \approx 1$ و شتاب بازگرداننده بیشینه است؛ در نزدیکی عمودی، $\sin\theta \to 0$ و شتاب به‌صورت طبیعی کاهش می‌یابد. هرچه $k_p^{\text{PFL}}$ بزرگ‌تر باشد، شتاب بازگرداننده قوی‌تر و همگرایی سریع‌تر است؛ اما مقادیر بیش از حد بزرگ، فرمان‌های شتاب چرخ بسیار بزرگی تولید می‌کنند که ممکن است از محدودیت‌های ولتاژ و سرعت چرخ فراتر روند.

\paragraph{نقش فیزیکی $k_d^{\text{PFL}}$ (بهره مشتقی).}
جمله $-k_d \dot\theta$ میرایی سرعتی فراهم می‌کند و نقش یک دمپر (میراگر) مجازی را ایفا می‌نماید. بدون این جمله، قانون مجازی به $\ddot\theta_{\text{des}} = -k_p \sin\theta$ تقلیل می‌یابد که یک نوسانگر بدون اتلاف است و پاندول به‌جای همگرایی به $\theta = 0$، در اطراف آن نوسان پایدار انجام می‌دهد. جمله میرایی، انرژی جنبشی مجازی سیستم را در هر نیم‌سیکل تخلیه می‌کند و نوسانات اطراف نقطه هدف را به‌تدریج از بین می‌برد. به بیان دیگر، $k_d^{\text{PFL}}$ مشخص می‌کند که پاندول با چه سرعتی از نوسان باز می‌ایستد و به حالت سکون در وضعیت عمودی می‌رسد.

\paragraph{تنظیم بهره‌ها.}
در نزدیکی وضعیت عمودی ($\sin\theta \approx \theta$)، قانون مجازی به معادله خطی مرتبه دوم زیر تبدیل می‌شود:

\begin{equation}
    \ddot\theta + k_d^{\text{PFL}}\, \dot\theta + k_p^{\text{PFL}}\, \theta = 0
\end{equation}

این معادله یک نوسانگر میرای مرتبه دوم است که دو پارامتر مشخصه دارد:

\begin{equation}
    \omega_n = \sqrt{k_p^{\text{PFL}}}, \qquad \zeta = \frac{k_d^{\text{PFL}}}{2\sqrt{k_p^{\text{PFL}}}}
\end{equation}

که در آن‌ها $\omega_n$ فرکانس طبیعی (تعیین‌کننده سرعت همگرایی) و $\zeta$ نسبت میرایی (تعیین‌کننده مقدار نوسان) است. بر این اساس، تنظیم بهره‌ها در دو گام انجام می‌شود:

\begin{enumerate}
    \item \textbf{انتخاب $\omega_n$:} فرکانس طبیعی مجازی باید به‌اندازه کافی بزرگ باشد که همگرایی در چند ثانیه حاصل شود، اما از پهنای باند عملی موتور فراتر نرود. محدوده مناسب $\omega_n \in [2,\; 3]$ \LR{rad/s} است. سپس $k_p^{\text{PFL}} = \omega_n^2$ قرار می‌دهیم.
    \item \textbf{انتخاب $\zeta$:} مقدار $\zeta = 1$ (میرایی بحرانی) سریع‌ترین همگرایی بدون نوسان را فراهم می‌کند. مقادیر $\zeta < 1$ پاسخ نوسانی میرا و $\zeta > 1$ همگرایی کندتر بدون نوسان تولید می‌کنند. سپس $k_d^{\text{PFL}} = 2\zeta\,\omega_n$ قرار می‌دهیم.
\end{enumerate}

\paragraph{قید شتاب و بررسی عددی.}
بهره‌های انتخابی باید در محدودیت فیزیکی موتور قابل تحقق باشند. حداکثر شتاب پاندول که موتور می‌تواند تولید کند، از روابط زیر به دست می‌آید:

\begin{equation}
    i_{\max} = \frac{V_{\max}}{R_a} = \frac{12}{1.2} = 10\;\text{A}, \qquad
    \tau_{\max} = K_t \, i_{\max} = 0.876\;\text{N·m}
\end{equation}

\begin{equation}
    \ddot\varphi_{\max} = \frac{\tau_{\max}}{I_{w,\text{eff}}} \approx 1230\;\text{rad/s}^2, \qquad
    \ddot\theta_{\max} \approx \ddot\varphi_{\max} \cdot \frac{M_{12}}{M_{11}} \approx 16\;\text{rad/s}^2
\end{equation}

در بدترین حالت ($\theta = \pi/2$ و $\dot\theta \approx 2$ \LR{rad/s})، شتاب مطلوب برابر $k_p^{\text{PFL}} + 2\,k_d^{\text{PFL}}$ است و باید کمتر از $16$ \LR{rad/s²} باقی بماند.

با انتخاب $\zeta = 1$ و $\omega_n = 2.2$ \LR{rad/s}:

\begin{equation}
    k_p^{\text{PFL}} = \omega_n^2 = 4.84 \approx 5.0, \qquad
    k_d^{\text{PFL}} = 2\zeta\,\omega_n = 4.4 \approx 4.5
\end{equation}

بررسی قید: $5.0 + 2 \times 4.5 = 14.0 < 16$ ✓. این مقادیر، میرایی بحرانی با زمان همگرایی حدود $2$ ثانیه فراهم می‌کنند و در محدوده عملی موتور قرار دارند. مقادیر پیش‌فرض پیاده‌سازی ($k_p^{\text{PFL}} = 5.0$، $k_d^{\text{PFL}} = 2.0$) متناظر با $\zeta \approx 0.45$ هستند و پاسخ نوسانی میرا تولید می‌کنند.

این دو جمله در کنار هم، یک میدان برداری پایدار به سمت $\theta = 0$ شکل می‌دهند. به بیان دیگر، اگر پاندول بتواند دقیقاً از $\ddot\theta_{\text{des}}$ پیروی کند، به‌صورت مجانبی به وضعیت عمودی همگرا می‌شود.

\subsubsection{معکوس‌سازی دینامیک}

اکنون پرسش اصلی این است: چه شتاب چرخی $\ddot\varphi$ لازم است تا پاندول شتاب $\ddot\theta_{\text{des}}$ را تجربه کند؟ برای پاسخ، از سطر اول معادلات کوپل‌شده حرکت شروع می‌کنیم:

\begin{equation}
    M_{11}\,\ddot\theta + M_{12}\,\ddot\varphi = G\sin\theta - N\,K_t\, i_a - b\,\dot\theta
\end{equation}

در این رابطه، سمت چپ شامل جملات اینرسی (شتاب‌های پاندول و چرخ) و سمت راست شامل نیروهای فعال (گرانش، گشتاور موتور و میرایی) است. حال $\ddot\theta = \ddot\theta_{\text{des}}$ را جایگذاری می‌کنیم و معادله را برای $\ddot\varphi$ حل می‌نماییم:

\begin{equation}
    M_{12}\,\ddot\varphi = G\sin\theta - N\,K_t\, i_a - b\,\dot\theta - M_{11}\,\ddot\theta_{\text{des}}
\end{equation}

و در نتیجه:

\begin{equation}
    \boxed{\ddot\varphi_{\text{req}} = \frac{G\sin\theta - N\,K_t\, i_a - b\,\dot\theta - M_{11}\,\ddot\theta_{\text{des}}}{M_{12}}}
\end{equation}

این رابطه، هسته خطی‌سازی بازخورد جزئی است. با محاسبه $\ddot\varphi_{\text{req}}$، عملاً جملات غیرخطی گرانش ($G\sin\theta$) و میرایی ($-b\dot\theta$) و اثر جریان لحظه‌ای آرمیچر ($-N K_t i_a$) را جبران می‌کنیم و رفتار پاندول را با رفتار خطی موردنظر ($\ddot\theta_{\text{des}}$) جایگزین می‌نماییم.

واژه «جزئی» در نام این روش به این دلیل است که فقط کانال پاندول خطی می‌شود. دینامیک داخلی چرخ (که از سطر دوم معادلات به دست می‌آید) همچنان غیرخطی باقی می‌ماند و می‌تواند مقادیر بزرگ $\ddot\varphi_{\text{req}}$ تولید کند. به همین دلیل، لایه‌های حفاظتی سرعت چرخ ضروری هستند.

نقش درایه غیرقطری $M_{12}$ در این رابطه محوری است. این کمیت، جفت‌شدگی سینتیکی بین پاندول و چرخ را بیان می‌کند و تنها مکانیزمی است که از طریق آن شتاب چرخ بر شتاب پاندول اثر می‌گذارد. با مقادیر پیش‌فرض، $M_{12}/M_{11} \approx 0.013$ است؛ بنابراین برای تولید یک واحد شتاب پاندول، به حدود $77$ واحد شتاب چرخ نیاز است. این نسبت، نیاز به سرعت‌های بالای چرخ در روش‌های بالاآوردن را توضیح می‌دهد.

\subsubsection{نگاشت گشتاور و ولتاژ}

پس از محاسبه $\ddot\varphi_{\text{req}}$، باید مشخص کنیم چه گشتاور موتوری این شتاب چرخ را هم‌زمان با حفظ $\ddot\theta_{\text{des}}$ تولید می‌کند. از سطر دوم معادلات کوپل‌شده استفاده می‌کنیم:

\begin{equation}
    M_{12}\,\ddot\theta + M_{22}\,\ddot\varphi = N\,K_t\, i_a - b_{w,\text{eff}}\,\dot\varphi
\end{equation}

با جایگذاری $\ddot\theta = \ddot\theta_{\text{des}}$ و $\ddot\varphi = \ddot\varphi_{\text{req}}$ و صرف‌نظر از جمله میرایی چرخ $b_{w,\text{eff}}\dot\varphi$ (که در مقایسه با جملات اینرسی کوچک است)، گشتاور موتور به صورت زیر به دست می‌آید:

\begin{equation}
    N\,K_t\, i_a = M_{12}\,\ddot\theta_{\text{des}} + M_{22}\,\ddot\varphi_{\text{req}}
\end{equation}

از آنجا که $\tau_m = K_t\, i_a$، گشتاور موتور مورد نیاز برابر است با:

\begin{equation}
    \boxed{\tau_m = \frac{M_{12}\,\ddot\theta_{\text{des}} + M_{22}\,\ddot\varphi_{\text{req}}}{N}}
\end{equation}

در اتصال مستقیم ($N = 1$)، این رابطه ساده‌تر می‌شود: $\tau_m = M_{12}\,\ddot\theta_{\text{des}} + M_{22}\,\ddot\varphi_{\text{req}}$.

گام نهایی، تبدیل گشتاور به ولتاژ است. مدار آرمیچر موتور DC از رابطه زیر پیروی می‌کند:

\begin{equation}
    V = R_a\, i_a + L_a \frac{di_a}{dt} + K_e\, N\, \dot\varphi
\end{equation}

ثابت زمانی الکتریکی $\tau_e = L_a / R_a \approx 0.42$ میلی‌ثانیه است، در حالی که مقیاس زمانی مکانیکی سیستم حدود $155$ میلی‌ثانیه (معادل $1/\omega_n$) است. این اختلاف دو مرتبه بزرگی، فرض شبه‌ایستای الکتریکی را توجیه می‌کند: $L_a\, di_a/dt \approx 0$. با این ساده‌سازی و جایگذاری $i_a = \tau_m / K_t$:

\begin{equation}
    \boxed{V = R_a \frac{\tau_m}{K_t} + K_e\, N\, \dot\varphi}
\end{equation}

هر یک از دو جمله این رابطه نقش فیزیکی مشخصی دارد:

\begin{itemize}
    \item \textbf{جمله اول} ($R_a\, \tau_m / K_t$): ولتاژ مقاومتی لازم برای ایجاد جریان آرمیچری که گشتاور $\tau_m$ را تولید کند. این جمله، بخش اصلی فرمان ولتاژ است.
    \item \textbf{جمله دوم} ($K_e\, N\, \dot\varphi$): جبران نیروی ضد محرکه الکتریکی (Back-EMF). چرخ در حال چرخش، ولتاژی متناسب با سرعت خود تولید می‌کند که با ولتاژ اعمالی مقابله می‌نماید. بدون این جبران، در سرعت‌های بالای چرخ، ولتاژ مؤثر آرمیچر به‌شدت کاهش می‌یابد و موتور قادر به تولید گشتاور مورد نیاز نخواهد بود.
\end{itemize}

\subsubsection{اشباع آگاه از انرژی}

اگر انرژی پاندول به هدف عمودی رسیده یا از آن فراتر رفته باشد ($E_p > 0$)، ادامه شتاب‌دهی باعث فراجهش از وضعیت عمودی می‌شود. در این حالت، کنترلر به حالت ترمز Back-EMF تغییر وضعیت می‌دهد:

\begin{equation}
    V = -K_e\, N\, \dot\varphi
\end{equation}

این فرمان معادل اتصال کوتاه الکتریکی موتور است: ولتاژ خارجی صفر قرار می‌گیرد و Back-EMF تولیدشده توسط چرخ، جریانی از طریق مقاومت آرمیچر هدایت می‌کند. این جریان، گشتاور ترمزی مخالف جهت چرخش ایجاد می‌کند و انرژی جنبشی چرخ را به حرارت در $R_a$ تبدیل می‌سازد. در نتیجه، سرعت چرخ کاهش می‌یابد و از شتاب‌دهی بیشتر پاندول جلوگیری می‌شود.

\subsection{بالاآوردن ضربه‌ای در سرعت صفر}

\subsubsection{شهود فیزیکی}

این روش مستقیم‌ترین راهبرد است: صبر می‌کنیم تا پاندول به نقطه بازگشت نوسان برسد (جایی که لحظه‌ای متوقف می‌شود)، سپس یک ضربه گشتاور شدید به چرخ وارد می‌کنیم. بر اساس پایستگی تکانه زاویه‌ای، پاندول یک ضربه زاویه‌ای مساوی و در جهت مخالف دریافت می‌کند و با انرژی بیشتری به مسیر خود بازمی‌گردد. تکرار این فرایند در هر نقطه بازگشت، دامنه نوسان را به‌تدریج افزایش می‌دهد.

\subsubsection{تشخیص عبور از صفر}

پاندول در نقطه بازگشت نوسان قرار دارد وقتی سرعت زاویه‌ای آن تغییر علامت دهد:

\begin{equation}
    (\dot\theta_{\text{prev}} > 0 \;\wedge\; \dot\theta \leq 0) \quad \lor \quad (\dot\theta_{\text{prev}} < 0 \;\wedge\; \dot\theta \geq 0)
\end{equation}

ضربه فقط در صورتی فعال می‌شود که $|\theta| > \theta_{\min} = 0.05$ رادیان باشد (برای جلوگیری از تحریک کاذب در نزدیکی تعادل عمودی) و انرژی پاندول هنوز کمتر از هدف باشد ($E_p < 0$).

\subsubsection{قانون ضربه}

در نقطه بازگشت شناسایی‌شده، یک ولتاژ ثابت به مدت $\Delta t_{\text{imp}}$ اعمال می‌شود:

\begin{equation}
    V = \operatorname{sign}(\theta) \cdot V_{\max}, \qquad t \in [0,\; \Delta t_{\text{imp}}]
\end{equation}

تعداد گام‌های انتگرال‌گیری $\lceil \Delta t_{\text{imp}} / h \rceil$ است که $h$ گام زمانی فیزیک است. قرارداد علامت $\operatorname{sign}(\theta)$ تضمین می‌کند که گشتاور واکنشی چرخ، پاندول را در نوسان برگشتی به سمت عمودی هل می‌دهد. وقتی $\theta > 0$ (پاندول به راست کج شده)، ولتاژ مثبت چرخ را در جهت مثبت شتاب می‌دهد و واکنش آن پاندول را به سمت عمودی بازمی‌گرداند.

\subsubsection{قطع ایمنی}

اگر در حین یک ضربه فعال، سرعت چرخ از $0.9\,\omega_{\max}$ فراتر رود، ضربه فوراً قطع شده و ولتاژ به صفر بازمی‌گردد. این مکانیزم از سرعت بیش از حد چرخ ناشی از ضربات متوالی جلوگیری می‌کند.

\subsection{تحلیل مقایسه‌ای و انتخاب روش بهینه}

پس از معرفی سه روش بالاآوردن، این پرسش مطرح می‌شود که کدام یک برای سیستم مورد بررسی مناسب‌تر است. پاسخ به این پرسش را می‌توان از دل خود دینامیک سیستم و با محاسبه نسبت‌های بی‌بعد مشخصه استخراج کرد. در این subsection، ابتدا نرخ انتقال انرژی از چرخ به پاندول را از معادلات کوپل‌شده مشتق می‌کنیم، سپس با ارزیابی عددی پارامترها، هر سه روش را در برابر محدودیت‌های فیزیکی سیستم می‌سنجیم.

\subsubsection{نسبت‌های مشخصه سیستم}

پیش از تحلیل، مقادیر عددی چند نسبت کلیدی را با پارامترهای پیش‌فرض جدول~\ref{tab:computed-values} محاسبه می‌کنیم:

\begin{table}[H]
\centering
\caption{نسبت‌های مشخصه تعیین‌کننده انتخاب الگوریتم بالاآوردن}
\label{tab:ratios}
\begin{tabular}{lccl}
\toprule
\textbf{کمیت} & \textbf{مقدار} & \textbf{مقیاس مکانیکی} & \textbf{نتیجه} \\
\midrule
ثابت زمانی الکتریکی $\tau_e = L_a/R_a$ & $0.42$ \LR{ms} & $T/2 \approx 490$ \LR{ms} & $\tau_e \ll T$ \\
نسبت کوپل $M_{12}/M_{11}$ & $0.013$ & --- & کوپل ضعیف \\
گشتاور بیشینه موتور $K_t V_{\max}/R_a$ & $0.876$ \LR{N·m} & $G = 2.21$ \LR{N·m} & $\tau_{\max} \approx 0.4\,G$ \\
تکانه زاویه‌ای چرخ در $\omega_{\max}$ & $0.036$ \LR{N·m·s} & $\mathcal{L}_p \approx 0.18$ \LR{N·m·s} & $\approx 20\%$ \\
\bottomrule
\end{tabular}
\end{table}

این اعداد چهار نتیجه فوری دارند:

\begin{enumerate}
    \item ثابت زمانی الکتریکی ($0.42$ میلی‌ثانیه) در مقایسه با نیم‌دوره نوسان پاندول ($490$ میلی‌ثانیه) سه مرتبه بزرگی کوچک‌تر است. بنابراین، جریان آرمیچر تقریباً بدون تأخیر از ولتاژ پیروی می‌کند و فرمان‌های ولتاژ بی‌درنگ به گشتاور تبدیل می‌شوند.
    \item نسبت کوپل $M_{12}/M_{11} \approx 0.013$ نشان می‌دهد که اینرسی چرخ تنها $1.3\%$ اینرسی کل سیستم است. برای تولید یک واحد شتاب پاندول، به حدود $77$ واحد شتاب چرخ نیاز است. این کوپل ضعیف، نیاز به چندین نوسان برای انباشت انرژی را اجتناب‌ناپذیر می‌سازد.
    \item گشتاور بیشینه موتور ($0.876$ \LR{N·m}) تنها $40\%$ گشتاور گرانشی ($2.21$ \LR{N·m}) است. موتور نمی‌تواند پاندول را به‌صورت استاتیکی در وضعیت افقی نگه دارد و صرفاً از طریق دینامیک ضربه‌ای و نوسانی قادر به بالاآوردن است.
    \item تکانه زاویه‌ای چرخ در سرعت بیشینه ($50$ \LR{rad/s}) حدود $20\%$ تکانه زاویه‌ای پاندول در اواسط نوسان است. این نسبت، تعداد نوسان‌های مورد نیاز برای بالاآوردن کامل را تعیین می‌کند.
\end{enumerate}

\subsubsection{مشتق نرخ انتقال انرژی}

برای درک مکانیزم انتقال انرژی از چرخ به پاندول، مشتق زمانی انرژی پاندول $E_p = \tfrac{1}{2}I_p\dot\theta^2 + G(\cos\theta - 1)$ را محاسبه می‌کنیم:

\begin{equation}
    \frac{dE_p}{dt} = I_p\,\dot\theta\,\ddot\theta - G\sin\theta\,\dot\theta
    = \dot\theta\!\left(I_p\,\ddot\theta - G\sin\theta\right)
\end{equation}

از سطر اول معادلات کوپل‌شده، شتاب پاندول برابر است با:

\begin{equation}
    \ddot\theta = \frac{G\sin\theta - NK_t\, i_a - b\,\dot\theta - M_{12}\,\ddot\varphi}{M_{11}}
\end{equation}

با جایگذاری در رابطه انرژی و ساده‌سازی، بخش قابل‌کنترل نرخ تغییرات انرژی (جملاتی که ولتاژ $V$ بر آن‌ها اثر می‌گذارد) به صورت زیر به دست می‌آید:

\begin{equation}
    \boxed{\frac{dE_p}{dt}\bigg|_{\text{ctrl}}
    = -\frac{I_p}{M_{11}}\,\dot\theta\!\left(NK_t\, i_a + M_{12}\,\ddot\varphi\right)}
\end{equation}

این رابطه، هسته فیزیکی انتخاب الگوریتم بهینه است. دو نتیجه مستقیم از آن استخراج می‌شود:

\begin{itemize}
    \item \textbf{شرط تزریق انرژی:} برای اینکه $dE_p/dt > 0$ باشد، عبارت داخل پرانتز باید خلاف علامت $\dot\theta$ باشد. به بیان دیگر، وقتی پاندول در جهت مثبت حرکت می‌کند ($\dot\theta > 0$)، گشتاور موتور و شتاب چرخ باید در جهتی باشند که انرژی به پاندول تزریق شود.
    \item \textbf{نقش کوپل ضعیف:} ضریب $I_p/M_{11} \approx 0.56$ نشان می‌دهد که تنها حدود $56\%$ انرژی جنبشی پاندول مستقیماً از طریق اینرسی پاندول قابل کنترل است. باقی‌مانده از طریق جملات کوپل $M_{12}$ منتقل می‌شود که ضریب کوچک $0.013$ دارد.
\end{itemize}

با توجه به اینکه $\tau_e \ll \Delta t$، جریان آرمیچر به‌صورت شبه‌ایستا از رابطه $i_a \approx (V - K_e N\dot\varphi)/R_a$ پیروی می‌کند. بنابراین، فرمان ولتاژ $V$ تقریباً بی‌درنگ به گشتاور موتور و در نتیجه به نرخ انتقال انرژی تبدیل می‌شود.

\subsubsection{ارزیابی روش انرژی (آستروم–فوروتو)}

قانون کنترل $V = -k_e(E_{\text{target}} - E_p)\dot\varphi$ را در رابطه نرخ انرژی جایگذاری می‌کنیم. شتاب چرخ $\ddot\varphi$ متناسب با گشتاور موتور و در نتیجه متناسب با $V$ است. نرخ انتقال انرژی حاصل:

\begin{equation}
    \frac{dE_p}{dt}\bigg|_{\text{energy}} \propto (E_{\text{target}} - E_p)\,\dot\varphi\,\dot\theta
\end{equation}

این عبارت در طول یک دوره کامل نوسان، به‌طور متوسط مثبت است. دلیل این امر، همبستگی فاز بین $\dot\varphi$ و $\dot\theta$ از طریق کوپل $M_{12}$ است. ویژگی‌های کلیدی این روش عبارتند از:

\begin{itemize}
    \item \textbf{زمان‌بندی خودکار بهره:} کسری انرژی $(E_{\text{target}} - E_p)$ به‌عنوان یک بهره متغیر عمل می‌کند. در ابتدای بالاآوردن (انرژی بسیار منفی)، پمپاژ قوی انجام می‌شود؛ در نزدیکی عمودی (انرژی نزدیک صفر)، پمپاژ به‌صورت طبیعی ملایم می‌گردد.
    \item \textbf{عدم نیاز به معکوس‌سازی مدل:} این روش از ساختار کامل ماتریس اینرسی استفاده نمی‌کند و تنها به علامت $\dot\varphi$ و مقدار عددی $E_p$ نیاز دارد. در نتیجه، به خطاهای پارامتری در $M_{12}$ و $G$ حساس نیست.
    \item \textbf{پایداری لیاپانوف:} تابع کاندید $V_p = \tfrac{1}{2}(E_p - E_{\text{target}})^2$ دارای مشتق $\dot{V}_p \leq 0$ است که همگرایی مجانبی انرژی به هدف را تضمین می‌کند.
\end{itemize}

\subsubsection{ارزیابی روش خطی‌سازی بازخورد جزئی}

روش \LR{PFL} شتاب چرخ مورد نیاز را از رابطه $\ddot\varphi_{\text{req}} = (G\sin\theta - M_{11}\ddot\theta_{\text{des}})/M_{12}$ محاسبه می‌کند. با $M_{12} = 0.0007125$ \LR{kg·m²}، تقسیم بر این مقدار کوچک، حساسیت بالایی به خطاهای پارامتری ایجاد می‌کند:

\begin{equation}
    \frac{\partial \ddot\varphi_{\text{req}}}{\partial M_{12}} = -\frac{G\sin\theta - M_{11}\ddot\theta_{\text{des}}}{M_{12}^2}
\end{equation}

یک خطای $10\%$ در $M_{12}$، خطای $10\%$ در $\ddot\varphi_{\text{req}}$ و در نتیجه خطای معادل در ولتاژ فرمان تولید می‌کند. علاوه بر این، فرض شبه‌ایستای الکتریکی ($L_a\, di_a/dt \approx 0$) در این روش نقش محوری دارد. نسبت $\tau_e / \Delta t = 0.42$ نشان می‌دهد که این فرض در فرکانس $1$ کیلوهرتز قابل‌قبول است، اما در گذراهای سریع (نزدیک اشباع ولتاژ)، تأخیر فاز کوچکی ایجاد می‌شود.

مزیت اصلی \LR{PFL}، همگرایی سریع‌تر در هر نوسان است. این روش به‌جای پمپاژ تدریجی انرژی، مستقیماً شتاب پاندول را فرمان می‌دهد و در صورت دقت مدل، مسیر بهینه‌تری را طی می‌کند.

\subsubsection{ارزیابی روش ضربه‌ای}

روش ضربه‌ای ولتاژ کامل را در لحظه عبور سرعت پاندول از صفر اعمال می‌کند. در این لحظه، $\dot\theta \approx 0$ برقرار است. با مراجعه به رابطه نرخ انتقال انرژی:

\begin{equation}
    \frac{dE_p}{dt}\bigg|_{\text{ctrl}} = -\frac{I_p}{M_{11}}\,\dot\theta\!\left(NK_t\, i_a + M_{12}\,\ddot\varphi\right) \approx 0
\end{equation}

در نقطه بازگشت، $\dot\theta \approx 0$ و بنابراین نرخ انتقال انرژی به پاندول تقریباً صفر است. ضربه ولتاژ در این لحظه، عمدتاً انرژی جنبشی چرخ را افزایش می‌دهد و تنها در نوسان بعدی (وقتی $\dot\theta \neq 0$ می‌شود) بخشی از این انرژی به پاندول منتقل می‌گردد. این تأخیر ذاتی در انتقال انرژی، روش ضربه‌ای را از نظر بهره‌وری در رتبه آخر قرار می‌دهد.

\subsubsection{نتیجه‌گیری و توصیه}

جدول~\ref{tab:comparison} مقایسه سه روش را بر اساس معیارهای استخراج‌شده از دینامیک سیستم خلاصه می‌کند.

\begin{table}[H]
\centering
\caption{مقایسه روش‌های بالاآوردن بر اساس دینامیک سیستم}
\label{tab:comparison}
\begin{tabular}{lccc}
\toprule
\textbf{معیار} & \textbf{انرژی} & \textbf{\LR{PFL}} & \textbf{ضربه‌ای} \\
\midrule
تعداد نوسان تا عمودی & متوسط & کم & زیاد \\
حساسیت به خطای پارامتری & کم & زیاد & کم \\
بهره‌وری انتقال انرژی & بالا & بالا & پایین \\
نیاز به دقت مدل & ندارد & دارد & ندارد \\
پیچیدگی پیاده‌سازی & کم & متوسط & کم \\
\bottomrule
\end{tabular}
\end{table}

با توجه به کوپل ضعیف ($M_{12}/M_{11} \approx 0.013$)، گشتاور محدود موتور ($\tau_{\max} \approx 0.4\,G$) و ضرورت انباشت انرژی در چندین نوسان، \textbf{روش انرژی (آستروم–فوروتو) مناسب‌ترین انتخاب برای این سیستم} است. دلایل این انتخاب مستقیماً از دینامیک سیستم پیروی می‌کنند:

\begin{enumerate}
    \item کوپل ضعیف، بالاآوردن تک‌نوسانه را غیرممکن می‌سازد. روش انرژی به‌طور طبیعی انرژی را در چندین نوسان انباشته می‌کند.
    \item محدودیت گشتاور موتور، اشباع ولتاژ را در هر نوسان اجتناب‌ناپذیر می‌سازد. زمان‌بندی خودکار بهره در روش انرژی، از اشباع بی‌مورد جلوگیری می‌کند.
    \item حساسیت پایین به خطاهای پارامتری، این روش را در برابر عدم‌قطعیت مدل (که در سیستم‌های واقعی اجتناب‌ناپذیر است) مقاوم می‌سازد.
\end{enumerate}

در صورت نیاز به همگرایی سریع‌تر، یک رویکرد ترکیبی قابل‌تصور است: استفاده از روش انرژی برای $|\theta| > 0.5$ رادیان و سوئیچ به \LR{PFL} در ناحیه $0.3 < |\theta| < 0.5$ رادیان. با این حال، برای پارامترهای پیش‌فرض این سیستم، روش انرژی به‌تنهایی عملکرد کافی و قابل‌اطمینان ارائه می‌دهد.

\subsection{انتقال به کنترل تعادلی}

کنترلرهای بالاآوردن فقط در فاصله زیاد از وضعیت عمودی مؤثر هستند. در نزدیکی نقطه تعادل، یک کنترلر تعادلی پایدارساز کنترل را در دست می‌گیرد. شرط انتقال، قرار گرفتن هم‌زمان زاویه و سرعت در آستانه‌های مشخص است:

\begin{equation}
    |\theta| < \theta_{\text{up}} \quad \wedge \quad |\dot\theta| < \dot\theta_{\text{up}}
\end{equation}

که مقادیر پیش‌فرض $\theta_{\text{up}} = 0.3$ رادیان و $\dot\theta_{\text{up}} = 1.0$ رادیان بر ثانیه هستند.

پس از برقراری شرط انتقال، بسته به حالت انتخابی، کنترلر LQR یا PID فعال می‌شود. برای فرمول‌بندی کامل این کنترل‌کننده‌ها، به بخش~\ref{sec:controllers} مراجعه شود.

\subsection{قیدهای ایمنی سراسری}

تمام روش‌های بالاآوردن در یک طرح حفاظتی چندلایه پیچیده شده‌اند که مستقل از الگوریتم انتخابی عمل می‌کند. جدول~\ref{tab:safety} این لایه‌ها را خلاصه می‌کند.

\begin{table}[H]
\centering
\caption{لایه‌های حفاظتی سرعت چرخ و ولتاژ}
\label{tab:safety}
\begin{tabular}{lll}
\toprule
\textbf{لایه} & \textbf{شرط} & \textbf{اثر} \\
\midrule
محافظ سخت سرعت & $|\dot\varphi| > 1.5\,\omega_{\max}$ & $V = 0$ (قطع فوری) \\
فرمانده سرعت & $|\dot\varphi| \geq \omega_{\max}$ & $V = 0$ \\
کاهش باند بالا & $|\dot\varphi| > 0.6\,\omega_{\max}$ & $V \leftarrow V \cdot \dfrac{\omega_{\max} - |\dot\varphi|}{0.4\,\omega_{\max}}$ \\
اشباع ولتاژ & همواره & $V \leftarrow \operatorname{clip}(V,\; -V_{\max},\; +V_{\max})$ \\
\bottomrule
\end{tabular}
\end{table}

این لایه‌ها تضمین می‌کنند که چرخ هرگز از محدودیت‌های مکانیکی فراتر نرود، صرف‌نظر از پرخاشگری کنترل‌کننده یا تنظیم نادرست پارامترها.

\subsection{خلاصه پارامترها}

جدول~\ref{tab:swingup-params} پارامترهای اصلی روش‌های بالاآوردن را فهرست می‌کند.

\begin{table}[H]
\centering
\caption{پارامترهای روش‌های بالاآوردن}
\label{tab:swingup-params}
\begin{tabular}{lccl}
\toprule
\textbf{نماد} & \textbf{نام پارامتر} & \textbf{پیش‌فرض} & \textbf{نقش} \\
\midrule
$k_e$ & \texttt{energy\_swing\_up\_gain} & $1.0$ & بهره تزریق انرژی \\
$k_p^{\text{PFL}}$ & \texttt{pfl\_kp} & $5.0$ & بهره تناسبی PFL \\
$k_d^{\text{PFL}}$ & \texttt{pfl\_kd} & $2.0$ & بهره مشتقی PFL \\
$\omega_{\max}$ & \texttt{swing\_up\_max\_wheel\_speed} & $50.0$ rad/s & سقف سرعت چرخ \\
$\Delta t_{\text{imp}}$ & \texttt{zero\_velocity\_impulse\_duration} & $0.05$ s & مدت ضربه \\
$\varphi_{\text{exc}}$ & (ثابت کلاس) & $0.5$ rad/s & آستانه تحریک \\
$\alpha$ & (ثابت کلاس) & $0.3$ & کسر تحریک از $V_{\max}$ \\
$\theta_{\text{up}}$ & \texttt{upright\_angle\_threshold} & $0.3$ rad & آستانه زاویه انتقال \\
$\dot\theta_{\text{up}}$ & \texttt{upright\_velocity\_threshold} & $1.0$ rad/s & آستانه سرعت انتقال \\
$\theta_{\min}$ & (ثابت کلاس) & $0.05$ rad & محافظ تشخیص عبور از صفر \\
\bottomrule
\end{tabular}
\end{table}

\subsection{انتخاب حالت}

حالت کنترل تعیین می‌کند که آیا انتقال به کنترل تعادلی فعال است یا خیر. جدول~\ref{tab:mode-select} ترکیب‌های ممکن را نشان می‌دهد.

\begin{table}[H]
\centering
\caption{ترکیب حالت‌های بالاآوردن و تعادلی}
\label{tab:mode-select}
\begin{tabular}{lcc}
\toprule
\textbf{حالت} & \textbf{بالاآوردن} & \textbf{تعادلی} \\
\midrule
\texttt{swing\_up} / \texttt{energy\_swing\_up} & \checkmark & --- \\
\texttt{swing\_up\_lqr} & \checkmark & LQR \\
\texttt{swing\_up\_pid} & \checkmark & PID \\
\bottomrule
\end{tabular}
\end{table}

روش بالاآوردن به‌صورت مستقل از طریق پارامتر \texttt{swing\_up\_method} $\in$ \{energy, pfl, zero\_velocity\} انتخاب می‌شود. هر روش را می‌توان با هر حالت تعادلی ترکیب کرد.